% Generated draft: requires mathematical and source review.
\documentclass[11pt]{article}
\usepackage[margin=25mm]{geometry}
\usepackage{fontspec}
\setmainfont{Latin Modern Roman}
\usepackage{amsmath,amssymb,mathtools,unicode-math}
\setmathfont{Latin Modern Math}
\usepackage{enumitem,needspace,xurl,hyperref}
\setlength{\parindent}{0pt}
\setlength{\parskip}{5pt}
\setlength{\emergencystretch}{3em}
\newcommand{\N}{\mathbb{N}}
\newcommand{\Z}{\mathbb{Z}}
\newcommand{\Q}{\mathbb{Q}}
\newcommand{\R}{\mathbb{R}}
\newcommand{\C}{\mathbb{C}}
\newcommand{\F}{\mathbb{F}}
\newcommand{\A}{\mathbb{A}}
\newcommand{\PP}{\mathbb{P}}
\newcommand{\E}{\mathbb{E}}
\newcommand{\eps}{\varepsilon}
\newcommand{\dd}{\,\mathrm{d}}
\begin{document}
\section*{UnsolvedMath ID: 1100301}
\textbf{Source problem code:} catalogue entry
\subsection{Definitions and mathematical statement}
For a free group $K$, write $\operatorname{rk}(K)$ for the cardinality of a free basis; in the finite-rank setting define its reduced rank by
\[
 \overline r(K)=\max\{\operatorname{rk}(K)-1,0\}.
\]
The trivial group has rank zero. The precise finite-rank Hanna Neumann assertion identified in the supplied assessment is: for every free group $F$ and every pair of nontrivial finitely generated subgroups $A,B\leq F$, the intersection is of finite rank and
\[
 \operatorname{rk}(A\cap B)-1
 \leq(\operatorname{rk}(A)-1)(\operatorname{rk}(B)-1).
\]
Equivalently, $\overline r(A\cap B)\leq\overline r(A)\overline r(B)$, including the case of trivial intersection.

The supplied assessment also reports the strengthened double-coset inequality. Its standard finite-rank formulation is
\[
 \sum_{AgB\in A\backslash F/B}
 \overline r(A\cap gBg^{-1})
 \leq \overline r(A)\overline r(B),
\]
where one representative $g$ is chosen for each double coset; the summands are independent of that choice. The nonnegative sum is understood over all double cosets, including zero terms.

The original transcribed question says only ``nontrivial subgroups'' and does not explicitly require finite generation. The assessment explicitly gives the finitely generated formulation above. Without a convention for infinite-rank subtraction and multiplication, the original unrestricted rank expression is not a single specified numerical assertion; no infinite-rank extension is silently supplied.
\subsection{Short English statement}
The Hanna Neumann question bounds the rank of the intersection of two finitely generated free subgroups by the product of their reduced ranks. The supplied assessment identifies it, and the stronger sum over conjugate intersections, as proved independently by Mineyev and Friedman, with a simplification by Dicks. The registry labels it solved; these resolution statements are attributed to the supplied literature assessment, not to a proof or new verification in this draft.
\subsection{Sources}
\begin{itemize}
\item UnsolvedMath permanent ID 1100301, source code catalogue entry, ``Questions in Geometric Group Theory --- Q 3.1.'' Registry curation and metadata: Ulam AI and UnsolvedMath Contributors, CC BY 4.0; updated 2026-08-06.
\item Bestvina, \emph{Questions in Geometric Group Theory} (source list dated 2004), Question 3.1, Hanna Neumann Conjecture, PDF page 10: \url{https://www.math.utah.edu/~bestvina/eprints/questions-updated.pdf}. Supplied manual transcription accessed 2026-07-29. The original-list metadata records partial progress or special cases and calls for current-status release review; rights review is NEEDS\_REVIEW.
\item The supplied resolving assessment cites Mineyev, \emph{Annals of Mathematics} (2012), and Friedman, \emph{Memoirs of the American Mathematical Society} (2015), with a simplification by Dicks. Article titles, volumes, pages, locators, and a separate Dicks citation are not provided in the record; no missing details are inferred.
\end{itemize}
\end{document}
